\documentclass[12pt]{article}

% Packages
\usepackage{amsmath,amssymb,amsfonts,amsthm}
\usepackage{geometry}
\usepackage{hyperref}
\usepackage{booktabs}
\usepackage{enumitem}
\usepackage{graphicx}
\usepackage{float}
\usepackage{xcolor}
\usepackage{mathrsfs}

\geometry{margin=1in}

% Theorem environments
\newtheorem{theorem}{Theorem}[section]
\newtheorem{lemma}[theorem]{Lemma}
\newtheorem{proposition}[theorem]{Proposition}
\newtheorem{corollary}[theorem]{Corollary}
\newtheorem{conjecture}[theorem]{Conjecture}
\theoremstyle{definition}
\newtheorem{definition}[theorem]{Definition}
\newtheorem{example}[theorem]{Example}
\newtheorem{axiom}{Axiom}
\theoremstyle{remark}
\newtheorem{remark}[theorem]{Remark}

% Custom commands
\newcommand{\C}{\mathcal{C}}               % Clinical decision complex
\newcommand{\M}{\mathcal{M}}               % General manifold
\newcommand{\R}{\mathbb{R}}
\newcommand{\Z}{\mathbb{Z}}
\newcommand{\CF}{\mathrm{CF}}              % Clinical friction
\newcommand{\BI}{\mathrm{BI}}              % Bond Index
\newcommand{\MI}{\mathrm{MI}}              % Moral injury index

\title{\textbf{Geometric Clinical Ethics:\\Medical Decision-Making as Pathfinding on the Moral Manifold,\\and the Mathematical Structure of Moral Injury}}

\author{}

\date{}

\begin{document}

\maketitle

\noindent\textbf{Keywords:} clinical ethics, moral injury, medical decision-making, triage, informed consent, QALY, resource allocation, A* pathfinding, geometric ethics, moral manifold

\begin{abstract}
The dominant quantitative framework for medical decision-making---the quality-adjusted life year (QALY)---collapses multi-dimensional clinical-moral reality to a single scalar.  We provide the mathematical framework that clinical ethics has lacked.  Drawing on the Geometric Ethics programme, we construct the \emph{clinical decision complex}: a nine-dimensional weighted simplicial complex on which clinical decisions are pathfinding problems.  The clinician's decision is modeled as $A^*$ search with $f(n) = g(n) + h(n)$, where $g(n)$ is the accumulated clinical cost (evidence-based medicine) and $h(n)$ is the moral-heuristic cost (trained ethical intuition).  The minimum-cost path is the \emph{clinical geodesic}.

Six results follow.  (1)~Beauchamp and Childress's four principles are dimensions of the clinical manifold, and conflicts between them are resolved by the manifold's metric, not by ad hoc balancing.  (2)~The QALY destroys irrecoverable moral information; QALY-optimal policies can be manifold-suboptimal.  (3)~Informed consent is a gauge-invariance condition: consent is valid only if the patient's decision survives reframing.  (4)~Moral injury is mathematically characterized as cumulative forced boundary crossing---structurally distinct from burnout.  (5)~Triage is constrained multi-agent pathfinding; the framework predicts when utilitarian and manifold-optimal triage diverge.  (6)~The clinical Bond Index quantifies the moral cost imposed on patients by institutional policies, enabling measurement of structural healthcare injustice.

The framework generates falsifiable predictions about moral distress, consent quality, and health equity that differ from predictions of standard bioethics and QALY-based health economics.
\end{abstract}

%==============================================================================
\section{Introduction}
\label{sec:intro}
%==============================================================================

A physician who prescribes chemotherapy is simultaneously making a clinical judgment (will this extend survival?), an ethical judgment (does the patient understand the risks?), a resource-allocation judgment (is this the best use of a scarce drug?), and a relational judgment (does the patient trust me?).  These are not separable judgments.  They are \emph{simultaneous dimensions of a single decision}.

Yet clinical ethics has no mathematics.  Principlism \cite{beauchamp2019} provides four principles (autonomy, beneficence, non-maleficence, justice) but no algorithm for resolving conflicts between them.  QALY-based health economics \cite{weinstein1977,drummond2015} reduces outcomes to a scalar---mathematically tractable, morally catastrophic.  Casuistry \cite{jonsen1988} reasons by analogy but cannot handle novel cases.  Narrative ethics \cite{charon2006} is humane but non-quantitative.

We provide the missing mathematics by porting the Geometric Ethics framework \cite{bond2026geometric}---validated across 20,030 texts, 11 languages, and independently replicated \cite{thiele2026}---into clinical medicine.  The clinician computes $f(n) = g(n) + h(n)$, where $g(n)$ is accumulated clinical evidence (evidence-based medicine) and $h(n)$ is the moral-heuristic estimate (clinical wisdom).  The minimum-cost path on the full manifold is the \emph{clinical geodesic}.


%==============================================================================
\section{Related Work}
\label{sec:related}
%==============================================================================

\textbf{Principlism.}  Beauchamp and Childress's four principles \cite{beauchamp2019}---autonomy, beneficence, non-maleficence, justice---have dominated clinical ethics since 1979.  As Clouser and Gert \cite{clouser1990} argued, the principles are ``chapter headings,'' not action guides: when principles conflict, principlism offers no resolution algorithm.  Our framework dissolves this: the four principles map onto manifold dimensions ($d_4$, $d_1$, $d_2 + \beta_{\text{harm}}$, $d_3$ respectively), and ``conflicts'' are resolved by pathfinding on the metric, not by ad hoc balancing.

\textbf{QALY-based health economics.}  The QALY \cite{weinstein1977,drummond2015} reduces clinical reality to a scalar.  NICE uses \pounds20,000--\pounds30,000 per QALY gained as the standard threshold \cite{nice2013}.  Harris \cite{harris1987} and Nord et al.\ \cite{nord1999} have criticized QALYs for age and disability discrimination.  Our framework provides the mathematical basis: the discrimination arises from dimensional collapse (Theorem~\ref{thm:irrecoverability}).  A scalar that discards 8 of 9 morally relevant dimensions will systematically disadvantage populations whose needs are concentrated on the discarded dimensions.

\textbf{Moral injury.}  Jameton \cite{jameton1984} introduced moral distress; Dean et al.\ \cite{dean2019} reframed it as moral injury, borrowing from military psychology \cite{litz2009}.  Epstein and Hamric \cite{epstein2009} documented its cumulative nature (``moral residue'').  Our framework provides the first mathematical characterization: burnout is $g(n)$-depletion (resource exhaustion); moral injury is $h(n)$-damage (forced boundary crossing that permanently alters the clinician's heuristic function).

\textbf{Informed consent and shared decision-making.}  The shared decision-making model \cite{charles1997,elwyn2012} frames consent as a collaborative process.  Our framework formalizes this as \emph{manifold alignment}: the patient's $d_9$ (understanding) and $d_4$ (voluntariness) must be calibrated so that the patient's clinical geodesic approximates the geodesic they would compute with perfect information.

\textbf{Triage.}  Modern triage systems \cite{gilboy2011,emanuel2020} are explicitly utilitarian.  COVID-19 exposed their moral limits: withdrawing ventilators from current patients to reallocate them caused severe moral injury \cite{greenberg2020} because the utilitarian calculus operates on $d_1$ alone, while withdrawal crosses boundaries on $d_2$ (rights), $d_5$ (trust), and $d_7$ (identity).


%==============================================================================
\section{The Clinical Decision Complex}
\label{sec:manifold}
%==============================================================================

\subsection{Dimensions of Clinical Decision-Making}

Clinical decisions are evaluated on nine dimensions---the same nine dimensions of the moral manifold \cite{bond2026geometric}, re-interpreted for clinical contexts:

\begin{enumerate}
    \item $d_1$: \textbf{Clinical Outcomes} --- survival, morbidity, functional status, symptom burden.  This is the dimension that evidence-based medicine measures and that QALYs attempt to capture.
    \item $d_2$: \textbf{Rights and Obligations} --- the patient's right to treatment, the clinician's duty of care, Hippocratic non-maleficence (``first, do no harm''), the right to refuse treatment.
    \item $d_3$: \textbf{Justice and Fairness} --- equitable access, distributive justice in resource allocation, non-discrimination, proportionality of treatment burden.
    \item $d_4$: \textbf{Autonomy} --- the patient's right to self-determination, voluntariness of consent, freedom from coercion, the right to make ``bad'' decisions.
    \item $d_5$: \textbf{Trust} --- the therapeutic relationship, fiduciary duty, confidentiality, institutional trustworthiness.
    \item $d_6$: \textbf{Social and Relational Impact} --- effects on family and caregivers, community health consequences, social determinants of health, the patient's social role and responsibilities.
    \item $d_7$: \textbf{Dignity and Identity} --- the patient's sense of self, bodily integrity, personal values, quality of dying, the clinician's professional identity.
    \item $d_8$: \textbf{Institutional Legitimacy} --- standard of care, clinical guidelines, regulatory compliance, institutional protocols, legal defensibility.
    \item $d_9$: \textbf{Epistemic Status} --- diagnostic certainty, prognostic accuracy, the patient's understanding, health literacy, information quality.
\end{enumerate}

\begin{remark}[Why Nine Dimensions, Not Four]
Beauchamp and Childress's four principles map onto dimensions $d_1$--$d_4$ (approximately).  But clinical reality activates all nine.  A physician deciding whether to intubate an elderly patient with advanced dementia is simultaneously evaluating: clinical outcome ($d_1$: will intubation extend meaningful life?), rights ($d_2$: does the advance directive prohibit this?), justice ($d_3$: is this the best use of an ICU bed?), autonomy ($d_4$: what would the patient have wanted?), trust ($d_5$: what did I promise the family?), social impact ($d_6$: how will this affect the family's grief?), dignity ($d_7$: is this consistent with a dignified death?), institutional protocol ($d_8$: what does the hospital ethics committee recommend?), and epistemic status ($d_9$: how certain am I of the prognosis?).  Reducing this to four principles discards five dimensions.  Reducing it to a QALY discards eight.
\end{remark}

\subsection{The Clinical Decision Complex: Formal Construction}

\begin{definition}[Clinical Decision Complex]
\label{def:clinical_complex}
The \emph{clinical decision complex} $\C$ is a weighted simplicial complex constructed as follows:
\begin{itemize}
    \item \textbf{Vertices (0-simplices):} Each vertex $v_i$ represents a \emph{clinical state}---a configuration of the patient's condition, the available interventions, the clinician's knowledge, and the institutional context.  The vertex carries an \emph{attribute vector} $\mathbf{a}(v_i) \in \R^9$, whose components are scores along the nine clinical-moral dimensions.
    \item \textbf{Edges (1-simplices):} An edge $(v_i, v_j)$ represents an available \emph{clinical action}---a treatment decision, a diagnostic test, a referral, a conversation, a decision to wait.  The edge carries a weight $w(v_i, v_j) \geq 0$ representing the \emph{total cost} of the action on the full clinical manifold.
    \item \textbf{Higher simplices:} A $k$-simplex $[v_0, \ldots, v_k]$ represents a \emph{care bundle}---a set of jointly administered interventions (e.g., a chemotherapy protocol: drug + antiemetic + hydration + counseling).
\end{itemize}
\end{definition}

\begin{definition}[Clinical Edge Weights]
\label{def:clinical_weights}
The weight of a clinical action $(v_i, v_j)$ on $\C$ is:
\begin{equation}
w(v_i, v_j) = \underbrace{\Delta \mathbf{a}^T\, \Sigma^{-1}\, \Delta \mathbf{a}}_{\text{Mahalanobis distance}} + \underbrace{\sum_k \beta_k \cdot \mathbf{1}[\text{clinical-moral boundary } k \text{ crossed}]}_{\text{boundary penalties}}
\label{eq:clinical_weight}
\end{equation}
where $\Delta \mathbf{a} = \mathbf{a}(v_j) - \mathbf{a}(v_i)$ is the change in the patient's attribute vector, $\Sigma$ is the $9 \times 9$ \emph{clinical-moral covariance matrix} estimated from observed clinical decision-making, and $\beta_k$ is the penalty for crossing clinical-moral boundary $k$.
\end{definition}

The covariance matrix $\Sigma$ captures the cross-dimensional dependencies that make clinical ethics hard:
\begin{itemize}
    \item $\Sigma_{15}$ (outcomes $\times$ trust): a treatment with uncertain outcomes ($d_1$) is more costly when the trust relationship is fragile ($d_5$).
    \item $\Sigma_{49}$ (autonomy $\times$ epistemic): respecting autonomy ($d_4$) requires that the patient has adequate understanding ($d_9$); when $d_9$ is low, the effective cost of invoking $d_4$ increases.
    \item $\Sigma_{37}$ (justice $\times$ dignity): a resource-allocation decision that is fair ($d_3$) but strips dignity ($d_7$) from the patient (e.g., means-testing for treatment) has high cross-dimensional cost.
\end{itemize}

\subsection{Clinical-Moral Boundaries}

\begin{definition}[Clinical-Moral Boundaries]
\label{def:clinical_boundaries}
A \emph{clinical-moral boundary} is a threshold on $\C$ where the clinical decision regime changes discontinuously:

\begin{itemize}
    \item \textbf{The harm boundary ($\beta_{\text{harm}}$):} ``First, do no harm'' (Hippocratic non-maleficence).  Actions that directly cause patient harm cross this boundary.  For most clinicians, $\beta_{\text{harm}}$ is very high but finite---clinicians accept iatrogenic harm when the alternative (untreated disease) is worse.  This is why surgery exists: the harm boundary is crossed, but the clinical geodesic's total cost is lower with surgery than without.

    \item \textbf{The futility boundary ($\beta_{\text{futile}}$):} Continuing aggressive treatment when the clinical outcome is hopeless.  Crossing this boundary (treating a patient who cannot benefit) imposes costs on $d_3$ (justice: wasting scarce resources), $d_7$ (dignity: subjecting the patient to suffering without benefit), and $d_5$ (trust: promising what medicine cannot deliver).

    \item \textbf{The consent boundary ($\beta_{\text{consent}}$):} Treating a competent patient without valid consent.  For most clinicians, $\beta_{\text{consent}} \approx \infty$ for non-emergency interventions.  The boundary drops to a finite value in genuine emergencies (implied consent) and for patients who lack decision-making capacity (surrogate consent).

    \item \textbf{The abandonment boundary ($\beta_{\text{abandon}}$):} Withdrawing care from a patient who is dependent on the clinician.  The fiduciary nature of the therapeutic relationship ($d_5$) makes abandonment one of the highest-cost boundary crossings in clinical ethics.

    \item \textbf{The sacred-value boundary ($\beta_{\text{sacred}} = \infty$):}  Certain clinical actions are absolutely prohibited regardless of consequences: performing non-consensual experimentation, participating in execution, deliberately hastening death for non-palliative purposes.  These have $\beta = \infty$---no clinical benefit can offset them.
\end{itemize}
\end{definition}


%==============================================================================
\section{Clinical Decision as $A^*$ Search}
\label{sec:search}
%==============================================================================

\subsection{The Clinical Pathfinding Model}

\begin{definition}[Clinical Decision as Pathfinding]
\label{def:clinical_pathfinding}
A clinical decision is a pathfinding problem on $\C$:
\begin{itemize}
    \item \textbf{Start node $v_0$:} The patient's current clinical-moral state (diagnosis, prognosis, current treatments, relationships, institutional context).
    \item \textbf{Goal region $G$:} The set of acceptable clinical-moral outcomes (not just ``survival'' but survival with adequate function, autonomy, dignity, and trust).
    \item \textbf{Path $\gamma$:} A sequence of clinical actions (treatments, tests, conversations, decisions to wait) that moves the patient from $v_0$ toward $G$.
    \item \textbf{Cost function:} $f(n) = g(n) + h(n)$, where:
    \begin{itemize}
        \item $g(n)$: the \emph{accumulated clinical cost} from $v_0$ to the current state $n$---the sum of edge weights along the clinical path so far.  This is the \textbf{evidence-based medicine} component: slow, deliberate, data-driven, based on clinical trials, Bayesian updating, and outcome data.
        \item $h(n)$: the \emph{moral-heuristic estimate} of the remaining cost from $n$ to the nearest acceptable outcome in $G$---the clinician's trained ethical intuition about ``how far we still have to go'' and ``what moral costs remain.''  This is the \textbf{clinical wisdom} component: fast, automatic, based on experience, training, and internalized professional norms.
    \end{itemize}
\end{itemize}
\end{definition}

\begin{center}
\begin{tabular}{ll}
\toprule
\textbf{$A^*$ Component} & \textbf{Clinical Decision Counterpart} \\
\midrule
Start node $v_0$ & Patient's presenting condition \\
Goal region $G$ & Acceptable clinical-moral outcome \\
Edge cost $w(v_i, v_j)$ & Full cost of clinical action (all 9 dimensions) \\
$g(n)$: accumulated cost & Evidence-based medicine (EBM) \\
$h(n)$: heuristic estimate & Clinical wisdom / ethical intuition \\
Boundary penalty $\beta_k$ & Professional ethical norm (``do no harm'') \\
Open list & Treatment options under consideration \\
Closed list & Options already evaluated and dismissed \\
Optimal path $\gamma^*$ & Clinical geodesic: the recommended plan of care \\
Path not found & Clinical-moral impasse / ethics consultation \\
\bottomrule
\end{tabular}
\end{center}

\begin{definition}[Clinical Geodesic]
\label{def:clinical_geodesic}
The \emph{clinical geodesic} $\gamma^*$ from a patient's initial state $v_0$ to a goal region $G$ on the clinical decision complex $\C$ is the minimum-cost path:
\begin{equation}
\gamma^* = \arg\min_\gamma \sum_{i=0}^{m-1} w(v_i, v_{i+1}) \quad \text{subject to } \gamma(0) = v_0,\; \gamma(\text{end}) \in G
\label{eq:clinical_geodesic}
\end{equation}
where $w$ is the full multi-dimensional clinical edge weight (Equation~\ref{eq:clinical_weight}).
\end{definition}

\subsection{Clinical Heuristics as Admissible Search Heuristics}

Clinical training produces moral heuristics that function as $A^*$ search heuristics:

\begin{theorem}[Clinical Heuristic Admissibility]
\label{thm:clinical_admissibility}
Let $h_C(n)$ be a \emph{clinical moral heuristic}---a function that assigns high cost to actions violating professional ethical norms and low cost to actions consistent with them:
\begin{equation}
h_C(n) = \sum_{k} \beta_k \cdot P(\text{clinical action from } n \text{ crosses boundary } k)
\end{equation}
If the boundary penalties $\beta_k$ are calibrated by clinical training and professional socialization to be lower bounds on the true moral cost of boundary crossing (i.e., $\beta_k \leq \beta_k^*$, where $\beta_k^*$ includes malpractice liability, professional censure, patient harm, and psychological distress), then $h_C$ is admissible and $A^*$ search finds the optimal clinical path.
\end{theorem}

\begin{proof}
Follows from the general Heuristic Truncation Theorem \cite{bond2026geometric}: $h_C(n) \leq h^*(n)$ because $\beta_k \leq \beta_k^*$ and all remaining cost components (Mahalanobis distance to goal) are non-negative.
\end{proof}

\begin{remark}[Clinical Training as Heuristic Calibration]
Medical education, clinical rotations, and professional socialization are, in this framework, a \emph{heuristic calibration process}: the trainee's $h(n)$ function is being tuned to approximate the boundary penalties and cost estimates that experienced clinicians have internalized.  An attending physician's ``clinical judgment'' is a well-calibrated $h(n)$.  A medical student's ``gut feeling'' is an uncalibrated $h(n)$.  The goal of clinical training is not to teach the student to compute $g(n)$ (that is the role of medical school didactics and evidence-based medicine training) but to calibrate $h(n)$: to develop the moral-clinical intuition that allows fast, accurate assessment of the remaining cost to acceptable outcomes.
\end{remark}

\begin{remark}[The $\varepsilon$-Admissibility Spectrum]
\label{rem:epsilon}
Admissibility requires $h_C(n) \leq h^*(n)$, but human heuristics can overestimate.  The framework accommodates this via $\varepsilon$-admissibility: if $h_C(n) \leq (1+\varepsilon)\,h^*(n)$, the returned path costs at most $(1+\varepsilon)$ times the optimum \cite{hart1968}.  Clinical heuristics form a four-category hierarchy: (i)~\emph{strictly admissible}---prohibitions whose true cost exceeds any estimate (e.g., non-consensual treatment); (ii)~\emph{$\varepsilon$-admissible}---the normal operating range of well-trained clinicians, producing near-optimal paths at negligible cost; (iii)~\emph{inadmissible}---heuristics inflated by moral injury, trauma, or defensive medicine, causing clinicians to avoid beneficial interventions (this is the clinical signature of $h(n)$ damage, distinct from burnout); (iv)~\emph{gauge-variant}---heuristics sensitive to framing rather than clinical state (omission bias, anchoring), which alone deserve the label ``cognitive bias.''  Ethics consultations and M\&M conferences function as \emph{heuristic recalibration}: they restore inadmissible $h(n)$ toward $\varepsilon$-admissibility by correcting inflated boundary penalties.
\end{remark}

\subsection{The Scalar Irrecoverability Theorem in Clinical Context}

\begin{theorem}[QALY Irrecoverability]
\label{thm:irrecoverability}
Let $Q: \R^9 \to \R$ be any function that maps a 9-dimensional clinical-moral attribute vector to a scalar QALY value.  Then:
\begin{enumerate}
    \item $Q$ is not injective: multiple clinically and morally distinct states map to the same QALY score.
    \item The information destroyed by $Q$ is irrecoverable: there exists no function $\psi: \R \to \R^9$ such that $\psi \circ Q = \mathrm{id}$.
    \item The clinical geodesic on $\C$ is in general different from the QALY-maximizing path, and the difference is not bounded by any function of $Q$ alone.
\end{enumerate}
\end{theorem}

\begin{proof}
Identical to the Scalar Irrecoverability Theorem for economic utility \cite{bond2026geometric}: $Q$ maps $\R^9$ to $\R$; by the rank-nullity theorem, the kernel of $dQ$ has dimension $\geq 8$ at every regular point; information in the kernel is irrecoverable by the data processing inequality.
\end{proof}

\begin{corollary}[QALY Discrimination Is Mathematical, Not Moral]
\label{cor:qaly_discrimination}
The well-documented biases of QALY-based allocation against the elderly, the disabled, and minority populations \cite{harris1987,nord1999} are not failures of implementation.  They are mathematical consequences of dimensional collapse: populations whose health needs are concentrated on non-outcome dimensions ($d_4$: autonomy for disabled patients, $d_5$: trust for minority patients, $d_7$: dignity for elderly patients) are systematically disadvantaged by any scalar projection that weights $d_1$ (outcomes) most heavily.
\end{corollary}

\begin{proposition}[When QALYs Are Adequate]
\label{prop:qaly_adequate}
QALY-based analysis is an adequate approximation when and only when:
\begin{enumerate}
    \item The clinical decision activates primarily dimension $d_1$ (clinical outcomes), with negligible activation of $d_2$--$d_9$.
    \item The patient population is homogeneous on dimensions $d_2$--$d_9$ (no systematic variation in autonomy, trust, dignity, or justice concerns).
    \item No moral boundaries are crossed by any available action.
\end{enumerate}
These conditions are approximately satisfied for routine cost-effectiveness comparisons of pharmacologically similar drugs for a well-defined condition in a homogeneous population.  They are violated in every case involving: end-of-life care, mental health, disability, resource rationing, vulnerable populations, or cross-cultural care.
\end{proposition}


%==============================================================================
\section{The Mathematical Theory of Moral Injury}
\label{sec:moral_injury}
%==============================================================================

\subsection{Moral Injury Is Not Burnout}

Burnout is characterized by emotional exhaustion, depersonalization, and reduced personal accomplishment \cite{maslach1981}.  It is a \emph{resource depletion} phenomenon: the clinician's cognitive and emotional resources for computing $g(n)$ and $h(n)$ are exhausted by overuse.

Moral injury is structurally different.  It is not resource depletion but \emph{manifold damage}: the clinician is forced to traverse paths that cross moral boundaries, and the cumulative boundary-crossing cost produces lasting harm to the clinician's moral-psychological architecture.

\begin{definition}[Moral Injury as Forced Boundary Crossing]
\label{def:moral_injury}
A clinician experiences \emph{moral injury} at time $t$ if the institutional clinical geodesic $\gamma_{\text{inst}}^*$ (the path mandated by institutional policy, resource constraints, or hierarchical authority) crosses a moral boundary $k$ that the clinician's personal moral manifold assigns a boundary penalty $\beta_k^{(\text{clin})} > \tau$, where $\tau$ is the clinician's threshold for tolerable moral cost.  The \emph{moral injury increment} at time $t$ is:
\begin{equation}
\Delta \MI(t) = \sum_k \max\bigl(0,\; \beta_k^{(\text{clin})} - \tau\bigr) \cdot \mathbf{1}[\text{boundary } k \text{ crossed at time } t]
\label{eq:moral_injury}
\end{equation}
The \emph{cumulative moral injury} is:
\begin{equation}
\MI(T) = \sum_{t=1}^{T} \Delta \MI(t)
\label{eq:cumulative_mi}
\end{equation}
\end{definition}

\begin{theorem}[Moral Injury Accumulation]
\label{thm:mi_accumulation}
Moral injury is:
\begin{enumerate}
    \item \textbf{Cumulative:} $\MI(T)$ is monotonically non-decreasing in $T$.  Each forced boundary crossing adds to the total; no subsequent event subtracts from it.
    \item \textbf{Irreversible above threshold:} Once $\MI(T) > \MI_{\text{crit}}$ (a clinician-specific critical threshold), the clinician's heuristic function $h(n)$ is permanently altered---boundary penalties shift, trust in institutional legitimacy ($d_8$) degrades, and professional identity ($d_7$) is damaged.  This manifests clinically as the moral residue that Epstein and Hamric \cite{epstein2009} describe.
    \item \textbf{Structurally distinct from burnout:} A clinician can have high $\MI$ with low burnout (the work is sustainable but morally damaging---e.g., a well-rested ICU physician forced to implement a triage protocol that violates their ethical commitments) or low $\MI$ with high burnout (the work is exhausting but morally acceptable---e.g., a rural physician working excessive hours in an ethically supportive environment).
\end{enumerate}
\end{theorem}

\begin{proof}
(1) follows from the definition: $\Delta \MI(t) \geq 0$ for all $t$, so $\MI(T)$ is a non-decreasing sum.  (2) follows from the manifold-damage interpretation: forced boundary crossing does not merely impose a transient cost but permanently alters the topology of the clinician's decision manifold.  Boundaries that have been crossed involuntarily become either more sensitive (hypervigilance---the clinician becomes risk-averse on the violated dimension) or less sensitive (moral numbing---the boundary penalty decreases through repeated violation).  Both alterations are forms of lasting damage to $h(n)$.  (3) follows from the structural distinction between $g(n)$-depletion (burnout: the computation degrades) and $h(n)$-damage (moral injury: the heuristic itself is altered).
\end{proof}

\subsection{The Moral Injury Index}

\begin{definition}[Moral Injury Index]
\label{def:mi_index}
The \emph{Moral Injury Index} of an institutional policy $P$ is the expected cumulative moral injury imposed on a clinician operating under $P$:
\begin{equation}
\MI(P) = \mathbb{E}\left[\sum_{t} \sum_k \max\bigl(0,\; \beta_k^{(\text{clin})} - \tau\bigr) \cdot \mathbf{1}[P \text{ requires crossing boundary } k \text{ at } t]\right]
\label{eq:mi_index}
\end{equation}
where the expectation is over the distribution of clinical scenarios encountered under normal operations.
\end{definition}

\begin{proposition}[Institutional Moral Injury Is Predictable]
\label{prop:mi_predictable}
The Moral Injury Index $\MI(P)$ can be estimated \emph{before} a policy is implemented, by:
\begin{enumerate}
    \item Surveying clinicians to estimate $\beta_k^{(\text{clin})}$ (the boundary penalties in the clinician population's moral manifold).
    \item Analyzing the policy $P$ to identify which clinical scenarios require boundary crossing.
    \item Computing the expected frequency of those scenarios from historical case-mix data.
\end{enumerate}
A policy with $\MI(P) > \MI_{\text{crit}}$ is predicted to produce moral injury in the clinician workforce.  This prediction is testable against moral distress scores (e.g., the Moral Distress Scale--Revised, MDS-R \cite{hamric2012}).
\end{proposition}

\subsection{COVID-19 as Natural Experiment}

The pandemic forced clinicians to cross multiple moral boundaries simultaneously: denying ventilators ($\beta_{\text{abandon}}$), deprioritizing elderly patients ($\beta_{\text{justice}}$), restricting family visitation ($\beta_{\text{dignity}}$), and deciding under extreme uncertainty (degraded $d_9$).

\begin{proposition}[COVID-19 Moral Injury Prediction]
\label{prop:covid_mi}
The framework predicts pandemic moral injury should be: (1)~higher among ICU clinicians than non-triage specialties; (2)~higher among clinicians who disagreed with the triage protocol; (3)~correlated with the \emph{number} of triage decisions, not just severity; (4)~structurally dissociable from burnout (high MI with low burnout in clinicians making few but morally costly decisions); and (5)~partially predictable from pre-pandemic $\beta_k$ profiles.  Predictions (1)--(4) are consistent with the empirical literature \cite{greenberg2020,rushton2018}.  Prediction (5) is novel and testable.
\end{proposition}


%==============================================================================
\section{Geometric Informed Consent}
\label{sec:consent}
%==============================================================================

\subsection{Consent as Gauge Invariance}

The standard legal test for valid informed consent requires: (a) disclosure of material risks, (b) patient comprehension, (c) voluntariness, and (d) decision-making capacity.  These four requirements are typically treated as a checklist: if all four boxes are checked, consent is valid.

Our framework provides a deeper characterization.

\begin{definition}[Informed Consent as Manifold Calibration]
\label{def:consent_calibration}
Valid informed consent requires that the patient's clinical decision complex $\C_{\text{patient}}$ be \emph{calibrated}---that the patient's attribute vectors, edge weights, and boundary penalties approximate the true clinical-moral costs closely enough that the patient's clinical geodesic is a good approximation of the geodesic they would compute with perfect information.  Formally:
\begin{equation}
\|w_{\text{patient}}(v_i, v_j) - w_{\text{true}}(v_i, v_j)\| < \epsilon \quad \text{for all relevant edges } (v_i, v_j)
\label{eq:consent_calibration}
\end{equation}
where $\epsilon$ is the tolerance for manifold miscalibration.
\end{definition}

\begin{theorem}[Consent as Gauge Invariance]
\label{thm:consent_gauge}
Valid informed consent is equivalent to a \emph{gauge-invariance condition}: the patient's clinical decision should be invariant under changes in the \emph{description} of the options, dependent only on the \emph{attribute vectors} of the options.  Formally, if $\tau$ is a meaning-preserving transformation of the clinical information (e.g., reframing ``90\% survival rate'' as ``10\% mortality rate''), then consent is valid only if:
\begin{equation}
\gamma^*_{\text{patient}}(\tau(x)) = \gamma^*_{\text{patient}}(x) \quad \text{for all meaning-preserving } \tau
\end{equation}
If the patient's decision changes under reframing ($\tau$), then either $d_9$ (understanding) or $d_4$ (voluntariness) is inadequately calibrated, and consent is not genuinely informed.
\end{theorem}

\begin{proof}
The Bond Invariance Principle \cite{bond2026geometric} requires that ethical evaluations be invariant under meaning-preserving transformations.  If the patient's decision varies under $\tau$, the heuristic function $h_{\text{patient}}(n)$ is gauge-variant---it depends on the framing rather than on the underlying clinical reality.  This implies that $d_9$ (the patient's understanding of the options) is not calibrated to the true attribute vectors but is instead sensitive to surface features of the presentation.  Equivalently, the patient is not making an informed decision; they are making a \emph{framed} decision.
\end{proof}

\begin{corollary}[The Consent Paradox in Low Health Literacy]
\label{cor:consent_paradox}
Patients with low health literacy can sign consent forms (satisfying the legal checklist) while having $d_9$ values so low that their manifold is grossly miscalibrated.  The consent is legally valid but ethically vacuous: the patient's clinical geodesic on their miscalibrated manifold may be radically different from the geodesic they would compute with adequate understanding.  The gauge-invariance test detects this: if the patient's decision would change under reframing, the consent is not genuinely informed regardless of what the form says.
\end{corollary}

\subsection{Shared Decision-Making as Manifold Alignment}

\begin{definition}[Manifold Alignment]
\label{def:manifold_alignment}
Shared decision-making between clinician and patient is \emph{manifold alignment}: the process of bringing the patient's clinical decision complex $\C_{\text{patient}}$ into sufficient agreement with the clinician's complex $\C_{\text{clin}}$ on the relevant edges and vertices that a jointly computed clinical geodesic is acceptable to both parties.  This requires:
\begin{enumerate}
    \item \textbf{Epistemic alignment ($d_9$):} The patient understands the clinical facts (diagnosis, prognosis, treatment options, risks) with sufficient accuracy that their $d_1$ attribute vectors approximate the clinician's.
    \item \textbf{Value alignment ($d_2$--$d_8$):} The clinician understands the patient's boundary penalties ($\beta_k^{(\text{patient})}$) well enough that the clinical geodesic respects the patient's moral manifold, not just the clinician's.
    \item \textbf{Goal alignment:} The patient and clinician agree on the goal region $G$---what counts as an acceptable outcome.
\end{enumerate}
\end{definition}

\begin{remark}[When Alignment Fails]
Manifold alignment failure has three distinct clinical presentations: (i)~\emph{epistemic misalignment} ($d_9$ gap)---the patient does not understand the situation; (ii)~\emph{value misalignment} ($d_2$--$d_8$ gap)---the clinician recommends a path that crosses a boundary the patient considers inviolable; (iii)~\emph{goal misalignment} ($G$ disagreement)---clinician and patient disagree about acceptable outcomes.  Each failure mode requires a different intervention; lumping them under ``poor communication'' conflates three structurally distinct problems.
\end{remark}


%==============================================================================
\section{Geometric Triage}
\label{sec:triage}
%==============================================================================

\subsection{Triage as Constrained Multi-Agent Pathfinding}

Triage is not a single-patient clinical decision.  It is a \emph{multi-agent constrained pathfinding problem}: the clinician must compute clinical geodesics for multiple patients simultaneously, subject to resource constraints that make it impossible to place all patients on their individually optimal paths.

\begin{definition}[Triage Problem]
\label{def:triage}
A \emph{triage problem} is a tuple $(N, \C, R, G)$ where:
\begin{itemize}
    \item $N = \{p_1, \ldots, p_n\}$ is a set of patients.
    \item $\C = \{\C_1, \ldots, \C_n\}$ is the set of clinical decision complexes (one per patient).
    \item $R$ is a resource constraint: a bound on the total number of clinical actions (ventilators, ICU beds, surgical hours, medications) available.
    \item $G = \{G_1, \ldots, G_n\}$ is the set of goal regions (acceptable outcomes for each patient).
\end{itemize}
The triage decision is a feasible allocation of resources that minimizes total clinical-moral friction:
\begin{equation}
(\gamma_1^*, \ldots, \gamma_n^*) = \arg\min_{\gamma_1, \ldots, \gamma_n} \sum_{i=1}^n \CF(\gamma_i) \quad \text{subject to } \sum_{i=1}^n r(\gamma_i) \leq R
\label{eq:triage}
\end{equation}
where $\CF(\gamma_i)$ is the clinical friction (total path cost) for patient $i$'s path, and $r(\gamma_i)$ is the resource cost of that path.
\end{definition}

\begin{remark}[Tractability]
General multi-agent pathfinding is NP-hard, but clinical triage involves bounded problem sizes ($\sim$20--50 ICU beds), few resource types, and small per-patient action spaces.  The framework does not propose replacing existing greedy triage heuristics (e.g., ESI \cite{gilboy2011}) with exact solvers; it proposes that the \emph{objective function} should be total manifold cost $\sum_i \CF(\gamma_i)$, not total scalar outcome $\sum_i d_1(\gamma_i)$.  The existing algorithms work with the manifold objective---the change is in the cost function, not the algorithm.  The primary value is \emph{offline policy design}: evaluating crisis standards of care against the manifold objective before the crisis.
\end{remark}

\subsection{Utilitarian Triage vs.\ Manifold-Optimal Triage}

\begin{theorem}[Divergence of Utilitarian and Manifold-Optimal Triage]
\label{thm:triage_divergence}
Let $\gamma_U^*$ be the utilitarian triage allocation (minimize total $d_1$ cost: maximize lives or life-years saved) and let $\gamma_\C^*$ be the manifold-optimal allocation (minimize total clinical friction on all nine dimensions).  Then:
\begin{enumerate}
    \item $\gamma_U^* = \gamma_\C^*$ when the triage decisions activate only $d_1$ and no moral boundaries are crossed.
    \item $\gamma_U^* \neq \gamma_\C^*$ when triage decisions cross moral boundaries or activate non-outcome dimensions.  Specifically:
    \begin{enumerate}
        \item Withdrawing a ventilator from a current patient to give it to a new patient with better prognosis is utilitarian-optimal but may be manifold-suboptimal, because the withdrawal crosses the abandonment boundary ($\beta_{\text{abandon}}$) and trust boundary ($\beta_{\text{trust}}$) for the current patient.
        \item Age-based deprioritization (younger patients first) is utilitarian-optimal (maximizes life-years) but manifold-suboptimal when it violates $d_3$ (justice: equal moral worth regardless of age) for agents whose manifold weights fairness heavily.
        \item Lottery allocation (random assignment when prognoses are similar) is utilitarian-equivalent (no outcome difference) but manifold-superior: it preserves $d_3$ (fairness: procedural justice) and $d_8$ (legitimacy: transparent, rule-based process) while avoiding boundary crossings on $d_7$ (dignity: the patient is not devalued by the allocation criterion).
    \end{enumerate}
\end{enumerate}
\end{theorem}

\begin{proof}
(1) follows from the Mahalanobis distance reducing to the $d_1$ component when all other dimensions have zero activation.  (2a--2c) follow from the edge-weight structure: boundary penalties on $d_2$, $d_3$, $d_5$, $d_7$, and $d_8$ add positive cost to the utilitarian-optimal allocation, making the total manifold cost exceed that of an alternative allocation that avoids the boundary crossings.
\end{proof}

\begin{remark}[Why Clinicians Resist Utilitarian Triage]
The widespread clinician resistance to utilitarian crisis standards of care during COVID-19 \cite{greenberg2020,emanuel2020} is not irrationality and not squeamishness.  It is the clinician's $h(n)$ function correctly identifying that the utilitarian triage path crosses moral boundaries whose cost exceeds the utilitarian benefit.  The clinicians are right---on the full manifold.  The utilitarian algorithms are right---on the scalar projection.  The conflict is not ethical disagreement.  It is dimensional mismatch.
\end{remark}


%==============================================================================
\section{The Clinical Bond Index}
\label{sec:bond_index}
%==============================================================================

\subsection{Measuring Structural Healthcare Injustice}

\begin{definition}[Clinical Bond Index]
\label{def:clinical_bi}
The \emph{clinical Bond Index} $\BI(P, S)$ of an institutional policy $P$ with respect to a patient population $S$ is the expected deviation between the clinical geodesic on the full manifold and the path mandated by the policy:
\begin{equation}
\BI(P, S) = \mathbb{E}_{s \in S}\left[\CF(\gamma_P^*(s)) - \CF(\gamma_\C^*(s))\right]
\label{eq:clinical_bi}
\end{equation}
where $\gamma_P^*(s)$ is the path mandated by policy $P$ for patient $s$, $\gamma_\C^*(s)$ is the clinical geodesic for patient $s$ on the full manifold, and the expectation is over the patient population.
\end{definition}

\begin{proposition}[Bond Index Detects Structural Injustice]
\label{prop:structural_injustice}
If $\BI(P, S_1) \gg \BI(P, S_2)$ for two patient populations $S_1$ and $S_2$, then policy $P$ imposes systematically higher moral cost on population $S_1$.  This is a quantitative, measurable definition of structural healthcare injustice:
\begin{enumerate}
    \item \textbf{Racial disparities:} If Black patients have systematically lower trust ($d_5$) due to historical institutional betrayal \cite{washington2006}, and the policy assumes uniform trust, the clinical geodesic for Black patients is more costly than the policy recognizes.  The Bond Index detects this: $\BI(P, S_{\text{Black}}) > \BI(P, S_{\text{White}})$.
    \item \textbf{Disability discrimination:} If QALY-based allocation assigns lower quality weights to disabled patients (underweighting $d_4$: autonomy and $d_7$: dignity relative to $d_1$: function), disabled patients bear systematically higher manifold cost.  The Bond Index detects this.
    \item \textbf{Age discrimination:} If age-based triage criteria discount elderly patients' outcomes (lower remaining QALYs), the policy underweights $d_3$ (equal moral worth) and $d_7$ (dignity in aging).  The Bond Index detects this.
\end{enumerate}
\end{proposition}

\begin{remark}[The Bond Index Is Not a Subjective Judgment]
The clinical Bond Index is computed from measurable quantities: edge weights estimated from clinical decision data, boundary penalties estimated from clinician surveys, and population-level covariance matrices estimated from behavioral data.  The claim ``policy $P$ is structurally unjust toward population $S$'' becomes an empirical claim with a specific numerical value, not a philosophical opinion.  This enables evidence-based health equity policy: measure $\BI(P, S)$ for all relevant populations, identify policies with high $\BI$, and redesign them to reduce manifold cost.
\end{remark}


%==============================================================================
\section{Worked Examples}
\label{sec:examples}
%==============================================================================

\subsection{Example 1: The Competent Patient Who Refuses Treatment}

\textbf{Problem:} A 45-year-old Jehovah's Witness with acute gastrointestinal bleeding needs a blood transfusion.  Without transfusion, the probability of death is approximately 40\%.  The patient, competent and informed, refuses transfusion on religious grounds.

\textbf{Principlism:} Autonomy (respect the refusal) conflicts with beneficence (save the life).  Principlism says ``balance'' them but provides no algorithm.

\textbf{QALY analysis:} Transfusion gains $\sim$20 QALYs (remaining life expectancy).  Cost per QALY is negligible.  QALY-optimal action: transfuse.

\textbf{Clinical geodesic analysis:}
\begin{enumerate}
    \item \textbf{Path A (transfuse against refusal):} $\Delta d_1 = +\text{large}$ (survival), but $\Delta d_4 = -\infty$ (violating a competent refusal is an absolute autonomy violation), $\Delta d_2 = -\infty$ (violating the legal right to refuse treatment), $\Delta d_5 = -\text{large}$ (destroying the trust relationship), $\Delta d_7 = -\text{large}$ (violating the patient's deepest identity).  Edge weight: $w_A = \infty$.
    \item \textbf{Path B (respect refusal, optimize alternatives):} $\Delta d_1 = -\text{moderate}$ (higher mortality risk), but $\Delta d_4 = 0$ (autonomy intact), $\Delta d_2 = 0$ (rights respected), $\Delta d_5 = +\text{moderate}$ (trust reinforced by respecting the refusal), $\Delta d_7 = 0$ (identity intact).  Edge weight: $w_B = (\text{mortality risk})^2 / \sigma_1^2$.
    \item The clinical geodesic is Path B.  It is not ``ethically suboptimal'' or a ``concession to autonomy.''  It is the minimum-cost path on the full manifold.  Path A has infinite cost because $\beta_{\text{consent}} = \infty$.
    \item The framework also computes: explore alternative paths (erythropoietin, cell salvage, volume expanders) that reduce $\Delta d_1$ while maintaining zero cost on $d_2$--$d_9$.  The clinical geodesic includes aggressive pursuit of transfusion alternatives---not as a compromise but as the manifold-optimal strategy.
\end{enumerate}

\subsection{Example 2: Withdrawing Life Support}

\textbf{Problem:} An 82-year-old patient on ventilation for 3 weeks with no neurological recovery, no advance directive, and a divided family.

\textbf{Clinical geodesic analysis:}  On the manifold, $\Delta d_1 \approx 0$ (negligible survival benefit from continued ventilation), but $d_6$ (family), $d_5$ (trust), and $d_7$ (dignity) are all heavily activated.  The clinical geodesic is \emph{not} an immediate binary decision but a \emph{path through time}: family meeting $\to$ values conversation $\to$ time-limited trial $\to$ transition to comfort care.  A utilitarian analysis jumps to ``withdraw---no benefit.''  The clinical geodesic recognizes that the \emph{path} to withdrawal has lower total manifold cost than the immediate decision, even though the endpoint is the same.  The geodesic optimizes process, not just outcome.

%==============================================================================
\section{Falsifiable Predictions}
\label{sec:predictions}
%==============================================================================

The framework generates predictions that are specific, testable, and distinct from predictions of standard bioethics, QALY-based health economics, and existing moral distress research:

\begin{enumerate}
    \item \textbf{Prediction 1 (Dimensional Moral Distress):} Clinician moral distress (measured by MDS-R \cite{hamric2012}) should correlate with the \emph{number of manifold dimensions} activated in the distressing situation, not just with the severity on any single dimension.  \emph{Specific test:} Present clinicians with vignettes matched for clinical severity ($d_1$) but varying in the number of moral dimensions engaged.  Predicted: distress increases monotonically with dimension count.  \emph{What would falsify:} No correlation between dimension count and distress.

    \item \textbf{Prediction 2 (Moral Injury vs.\ Burnout Dissociation):} Moral injury (measured by the Moral Injury Symptom Scale, MISS-HP \cite{zhizhong2019}) and burnout (measured by the Maslach Burnout Inventory, MBI \cite{maslach1981}) should be statistically dissociable: high MI with low burnout should occur in clinicians facing frequent boundary-crossing decisions with adequate staffing, and low MI with high burnout should occur in clinicians facing high volume without boundary-crossing decisions.  \emph{What would falsify:} Perfect correlation between MI and burnout (they are not distinct constructs).

    \item \textbf{Prediction 3 (Consent Gauge-Invariance Test):} Patients with genuinely informed consent should show framing invariance: their treatment preferences should not change when clinically equivalent information is reframed (e.g., ``90\% survival'' vs.\ ``10\% mortality'').  Patients who show framing effects have inadequate $d_9$ calibration regardless of whether they signed the consent form.  \emph{Specific test:} Administer framing-variant and framing-invariant versions of consent information.  Measure treatment choice consistency.  Predicted: framing-variant patients have lower health literacy scores and lower self-reported understanding.  \emph{What would falsify:} No correlation between framing invariance and comprehension.

    \item \textbf{Prediction 4 (Bond Index Predicts Disparities):} The clinical Bond Index $\BI(P, S)$ computed from institutional policy analysis should predict patient-reported experience disparities across demographic groups.  \emph{Specific test:} Compute $\BI$ for a hospital's triage, consent, and discharge policies.  Correlate with HCAHPS patient experience scores stratified by race, age, and disability status.  Predicted: groups with higher $\BI$ report worse experience on trust, communication, and respect dimensions.  \emph{What would falsify:} No correlation between $\BI$ and patient experience disparities.

    \item \textbf{Prediction 5 (Triage Resistance Correlates):} Clinician resistance to utilitarian triage protocols should correlate with pre-crisis $\beta_k$ profiles.  \emph{What would falsify:} No correlation between $\beta_k$ profiles and triage resistance.

    \item \textbf{Prediction 6 (Process-Path Superiority):} In end-of-life decisions, satisfaction should be higher when the clinical geodesic includes process steps (family meetings, values conversations) than when the same endpoint is reached directly, \emph{even though the clinical outcome is identical}.  \emph{What would falsify:} No satisfaction difference between process-path and direct-path groups.
\end{enumerate}

The framework would be falsified by: evidence that moral distress is purely a function of $d_1$ severity; evidence that moral injury and burnout are empirically indistinguishable; evidence that QALY allocation produces no demographic disparities; or discovery of a clinical dilemma that cannot be decomposed into the nine-dimensional attribute space.


%==============================================================================
\section{Discussion}
\label{sec:discussion}
%==============================================================================

\subsection{Implications for Practice}

If the framework is correct, five consequences follow.  (1)~\emph{The QALY is supplemented, not abandoned}: QALY analysis remains valid for decisions that activate primarily $d_1$ (Proposition~\ref{prop:qaly_adequate}); for all others, health technology assessment should report multi-dimensional impact profiles.  (2)~\emph{Moral injury becomes predictable}: the Moral Injury Index enables prospective evaluation of policies before they damage clinicians.  (3)~\emph{Informed consent is testable}: does the patient's decision survive reframing?  If not, consent is procedurally but not ethically valid.  (4)~\emph{Health equity becomes measurable}: the Bond Index provides a numerical comparison of structural injustice across policies and populations.  (5)~\emph{Ethics consultation becomes computation}: manifold analysis supplements narrative deliberation by making the structure of the judgment explicit.

\subsection{Limitations and Open Problems}

The $9 \times 9$ covariance matrix $\Sigma$ has not been estimated from clinical data---the central empirical challenge.  We propose a concrete design: (i)~retrospective coding of $\geq 500$ ethics committee consultations along all nine dimensions to estimate the sample covariance, (ii)~a prospective discrete choice experiment using fractional-factorial clinical vignettes ($n \geq 200$ clinicians) to estimate marginal rates of substitution via conditional logit, and (iii)~cross-validation against held-out consult recommendations.  This requires no new instruments---only standard conjoint methodology applied to the manifold's dimensional structure.

For the Bond Index, six of nine dimensions are routinely measured in clinical quality registries ($d_1$, $d_2$, $d_6$, $d_8$).  The remaining three require validated proxies: $d_5$ (Trust) via the Wake Forest Trust in Physician Scale or HCAHPS communication composite; $d_7$ (Dignity) via the Patient Dignity Inventory \cite{chochinov2008} or the HCAHPS ``treated with respect'' item; $d_9$ (Epistemic Status) via health literacy instruments (REALM, NVS) combined with the framing-invariance test from Prediction~3.

The boundary penalties $\beta_k$ vary across clinical cultures (e.g., individual-autonomy vs.\ family-centered medicine); this variation must be empirically mapped.  The framework is a theoretical structure, not a bedside calculator; it informs decision support tools and policy analysis but is not intended for real-time clinical computation.


%==============================================================================
\section{Conclusion}
\label{sec:conclusion}
%==============================================================================

Clinical medicine has operated for decades with a gap between its ethical sophistication and its mathematical tools.  Clinicians reason on a multi-dimensional moral manifold every day---weighing outcomes against autonomy, justice against efficiency, dignity against survival.  The quantitative frameworks available to them (QALYs, utilitarian triage algorithms, cost-effectiveness ratios) operate on a scalar projection that discards most of the moral information they are trying to process.

This paper closes the gap.  The clinical decision complex $\C$ provides the formal structure.  The $A^*$ pathfinding model provides the computational architecture.  The clinical geodesic provides the solution concept.  The Moral Injury Index provides the workforce protection tool.  The clinical Bond Index provides the health equity measurement tool.

Clinical ethics is not a soft science that resists quantification.  It is a geometric science whose geometry has not, until now, been written down.

The manifold exists.  The clinicians are already pathfinding on it.  Now there is a mathematics for what they do.


%==============================================================================
% References
%==============================================================================
\begin{thebibliography}{99}

\bibitem{bond2026geometric}
A.~H. Bond, \emph{Geometric Ethics: The Mathematical Structure of Moral Reasoning}, v1.15, Department of Computer Engineering, San Jos\'{e} State University, Spring 2026.

\bibitem{thiele2026}
L.~Thiele, ``Independent Replication of MoralVector Decodability and Cross-Lingual Transfer in LaBSE Embeddings,'' Dept.\ Computer Science, UCLA, Tech.\ Rep., 2026.

\bibitem{beauchamp2019}
T.~L. Beauchamp and J.~F. Childress, \emph{Principles of Biomedical Ethics}, 8th ed.  Oxford University Press, 2019.

\bibitem{weinstein1977}
M.~C. Weinstein and W.~B. Stason, ``Foundations of Cost-Effectiveness Analysis for Health and Medical Practices,'' \emph{New England Journal of Medicine}, vol.~296, no.~13, pp.~716--721, 1977.

\bibitem{drummond2015}
M.~F. Drummond, M.~J. Sculpher, K.~Claxton, G.~L. Stoddart, and G.~W. Torrance, \emph{Methods for the Economic Evaluation of Health Care Programmes}, 4th ed.  Oxford University Press, 2015.

\bibitem{jonsen1988}
A.~R. Jonsen and S.~Toulmin, \emph{The Abuse of Casuistry: A History of Moral Reasoning}.  University of California Press, 1988.

\bibitem{charon2006}
R.~Charon, \emph{Narrative Medicine: Honoring the Stories of Illness}.  Oxford University Press, 2006.

\bibitem{clouser1990}
K.~D. Clouser and B.~Gert, ``A Critique of Principlism,'' \emph{The Journal of Medicine and Philosophy}, vol.~15, no.~2, pp.~219--236, 1990.

\bibitem{nice2013}
National Institute for Health and Care Excellence, ``Guide to the Methods of Technology Appraisal 2013,'' NICE Process and Methods Guides, London, 2013.

\bibitem{harris1987}
J.~Harris, ``QALYfying the Value of Life,'' \emph{Journal of Medical Ethics}, vol.~13, no.~3, pp.~117--123, 1987.

\bibitem{nord1999}
E.~Nord, J.~L. Pinto, J.~Richardson, P.~Menzel, and P.~Ubel, ``Incorporating Societal Concerns for Fairness in Numerical Valuations of Health Programmes,'' \emph{Health Economics}, vol.~8, no.~1, pp.~25--39, 1999.

\bibitem{jameton1984}
A.~Jameton, \emph{Nursing Practice: The Ethical Issues}.  Prentice-Hall, 1984.

\bibitem{epstein2009}
E.~G. Epstein and A.~B. Hamric, ``Moral Distress, Moral Residue, and the Crescendo Effect,'' \emph{The Journal of Clinical Ethics}, vol.~20, no.~4, pp.~330--342, 2009.

\bibitem{rushton2018}
C.~H. Rushton, \emph{Moral Resilience: Transforming Moral Suffering in Healthcare}.  Oxford University Press, 2018.

\bibitem{dean2019}
W.~Dean, S.~Talbot, and A.~Dean, ``Reframing Clinician Distress: Moral Injury Not Burnout,'' \emph{Federal Practitioner}, vol.~36, no.~9, pp.~400--402, 2019.

\bibitem{litz2009}
B.~T. Litz, N.~Stein, E.~Delaney, L.~Lebowitz, W.~P. Nash, C.~Silva, and S.~Maguen, ``Moral Injury and Moral Repair in War Veterans: A Preliminary Model and Intervention Strategy,'' \emph{Clinical Psychology Review}, vol.~29, no.~8, pp.~695--706, 2009.

\bibitem{charles1997}
C.~Charles, A.~Gafni, and T.~Whelan, ``Shared Decision-Making in the Medical Encounter: What Does It Mean? (or It Takes at Least Two to Tango),'' \emph{Social Science \& Medicine}, vol.~44, no.~5, pp.~681--692, 1997.

\bibitem{elwyn2012}
G.~Elwyn, D.~Frosch, R.~Thomson, N.~Joseph-Williams, A.~Lloyd, P.~Kinnersley, E.~Cording, D.~Tomson, C.~Dodd, S.~Rollnick, A.~Edwards, and M.~Barry, ``Shared Decision Making: A Model for Clinical Practice,'' \emph{Journal of General Internal Medicine}, vol.~27, no.~10, pp.~1361--1367, 2012.

\bibitem{gilboy2011}
N.~Gilboy, T.~Tanabe, D.~Travers, and A.~M. Rosenau, ``Emergency Severity Index (ESI): A Triage Tool for Emergency Department Care,'' Version~4, AHRQ Publication No.\ 12-0014, 2011.

\bibitem{emanuel2020}
E.~J. Emanuel, G.~Persad, R.~Upshur, B.~Thome, M.~Parker, A.~Glickman, C.~Zhang, C.~Boyle, M.~Smith, and J.~P. Phillips, ``Fair Allocation of Scarce Medical Resources in the Time of Covid-19,'' \emph{New England Journal of Medicine}, vol.~382, no.~21, pp.~2049--2055, 2020.

\bibitem{greenberg2020}
N.~Greenberg, M.~Docherty, S.~Gnanapragasam, and S.~Wessely, ``Managing Mental Health Challenges Faced by Healthcare Workers During Covid-19 Pandemic,'' \emph{BMJ}, vol.~368, m1211, 2020.

\bibitem{maslach1981}
C.~Maslach and S.~E. Jackson, ``The Measurement of Experienced Burnout,'' \emph{Journal of Organizational Behavior}, vol.~2, no.~2, pp.~99--113, 1981.

\bibitem{hamric2012}
A.~B. Hamric, C.~T. Borchers, and E.~G. Epstein, ``Development and Testing of an Instrument to Measure Moral Distress in Healthcare Professionals,'' \emph{AJOB Primary Research}, vol.~3, no.~2, pp.~1--9, 2012.

\bibitem{zhizhong2019}
W.~Zhizhong, J.~Koenig, S.~G.~Y. Yan, C.~Al Shohaib, and H.~G. Koenig, ``Psychometric Properties of the Moral Injury Symptom Scale Among Chinese Health Professionals During the COVID-19 Pandemic,'' \emph{Frontiers in Psychiatry}, vol.~11, 571354, 2020.

\bibitem{washington2006}
H.~A. Washington, \emph{Medical Apartheid: The Dark History of Medical Experimentation on Black Americans from Colonial Times to the Present}.  Doubleday, 2006.

\bibitem{hart1968}
P.~E. Hart, N.~J. Nilsson, and B.~Raphael, ``A Formal Basis for the Heuristic Determination of Minimum Cost Paths,'' \emph{IEEE Transactions on Systems Science and Cybernetics}, vol.~4, no.~2, pp.~100--107, 1968.

\bibitem{chochinov2008}
H.~M. Chochinov, T.~Hassard, S.~McClement, T.~Hack, L.~J. Kristjanson, M.~Harlos, S.~Sinclair, and A.~Murray, ``The Patient Dignity Inventory: A Novel Way of Measuring Dignity-Related Distress in Palliative Care,'' \emph{Journal of Pain and Symptom Management}, vol.~36, no.~6, pp.~559--571, 2008.

\end{thebibliography}

\end{document}
